% Generálva: html/build_tex.js (o06.js). Kézzel ne szerkeszd, a forrás a html/data mappa.
\setcounter{chapter}{5}
\chapter{\texorpdfstring{Forráskódolás és hibavédő kódolás}{Forráskódolás és hibavédő kódolás}}\label{f6}
Paritásbitekből generátormátrix, szindróma és javítás. Huffman-kód és entrópia, Kraft-egyenlőtlenség, és egy szokatlan feladat az állapotok kódolásáról.

\section{Alapok, egyszerűen}
\subsection{\texorpdfstring{A1. Jelölések: blokk-kód}{A1. Jelölések: blokk-kód}}\label{b6-jel}
\par\noindent{\small\renewcommand{\arraystretch}{1.25}\begin{tabularx}{\linewidth}{@{}>{\raggedright\arraybackslash}X>{\raggedright\arraybackslash}X@{}}\toprule
\textbf{Jel} & \textbf{Jelentés}\\\midrule
undefined & üzenet (undefined bit)\\
undefined & kódszó (undefined bit), undefined\\
undefined & vett szó, undefined, ahol undefined a hibavektor (1-es ott, ahol a bit elromlott)\\
undefined & generátormátrix (undefined), sorai az egységvektor-üzenetek kódszavai\\
undefined, undefined & paritásellenőrző mátrix és transzponáltja, undefined minden kódszóra\\
undefined & szindróma, undefined\\
undefined & kódtávolság: két kódszó közti legkisebb eltérés, lineáris kódnál a legkisebb nem nulla kódszósúly\\
\bottomrule\end{tabularx}}\par

\textbf{Recept.} Szisztematikus kódnál $G=[I\,|\,B]$ (elöl az üzenet, utána a paritásbitek), $H^T=\begin{bmatrix}B\\I\end{bmatrix}$, $H=[B^T\,|\,I]$. A szindróma a $v$ egyeseinek helyén álló $H^T$-sorok XOR-ja. Ha $s$ egyenlő $H^T$ valamelyik sorával, egyetlen hibát feltételezve az a bit romlott el. Biztosan jelezhető $d_{\min}-1$, javítható $\lfloor(d_{\min}-1)/2\rfloor$ hiba.

\subsection{\texorpdfstring{A2. Forráskódolás}{A2. Forráskódolás}}\label{b6-forras}
$p_i$ az $i$-edik szimbólum valószínűsége, $l_i$ a kódszavának hossza. Átlagos kódszóhossz $\bar L=\sum p_il_i$, entrópia $H=\sum p_i\log_2(1/p_i)$, és minden egyértelműen dekódolható kódra $\bar L\ge H$. \textbf{Kraft és McMillan}: egyértelműen dekódolható (és akkor prefix) kód pontosan akkor létezik a megadott hosszakkal, ha $\sum2^{-l_i}\le1$. \textbf{Huffman}: mindig a két legkisebb valószínűséget vonjuk össze, amíg egy marad. Ha minden $p_i$ kettő hatványa, $\bar L=H$.

% ------------------------------------------------------------------
\section{\texorpdfstring{6.1. feladat: Két kód paritásképletekből: C$_1$ és C$_2$}{6.1. feladat: Két kód paritásképletekből: C₁ és C₂}}\label{o6_01}
\setcounter{lepes}{0}
\begin{feladat}{6.1. feladat}
$C_1$: $u_1u_2u_3u_4\to u_1u_2u_3u_4p_1p_2p_3$, $p_1=u_1\oplus u_2\oplus u_4$, $p_2=u_1\oplus u_3\oplus u_4$, $p_3=u_2\oplus u_3\oplus u_4$. $C_2$: ugyanez és $p_4=p_1\oplus p_2\oplus p_3$. a) $C_1$ kódszavai. b) $C_1$ paritásellenőrző mátrixa, hány bites a szindróma? c) Mi lehetett az üzenet 0110000 vételekor, biztos? d) $C_2$ generátormátrixa?

{\footnotesize\sffamily Forrás: 5-ösöm 17\ldots{}18. oldal $\cdot$ 5o\_17\_1. Címkék: vizsga}
\end{feladat}
\begin{keresunk}
\begin{itemize}
\item a) a 16 kódszó
\item b) $H$ (paritásellenőrző mátrix), $H^T$ és a szindróma hossza
\item c) $\hat u$ (a becsült üzenet), és hogy biztos-e
\item d) $G_2$ (a második kód generátormátrixa)
\end{itemize}
\end{keresunk}
\begin{jeloles}
\begin{jeltabla}
$C_1$ & az első kód (7 bites kódszavak) & \\
$\ca{n}$ & a kódszó hossza bitben & $7$\\
$\cb{k}$ & az üzenet hossza bitben & $4$\\
$\cc{v}$ & a vett szó & $0110000$\\
$H^T$ & a paritásellenőrző mátrix transzponáltja & \emph{keresett}\\
$\hat u$ & a becsült üzenet & \emph{keresett}\\
$G_2$ & a második kód generátormátrixa & \emph{keresett}\\
$u_1$ & az üzenet 1. bitje & \\
$u_2$ & az üzenet 2. bitje & \\
$u_3$ & az üzenet 3. bitje & \\
$u_4$ & az üzenet 4. bitje & \\
$p_1$ & az 1. paritásbit & \\
$p_2$ & a 2. paritásbit & \\
$p_3$ & a 3. paritásbit & \\
$C_2$ & a második kód (8 bites kódszavak) & \\
$p_4$ & a 4. paritásbit (csak a $C_2$ kódban) & \\
$G$ & generátormátrix & \\
$u$ & az üzenet (4 bit) & \\
$I$ & egységmátrix & \\
$B$ & a generátormátrix paritásrésze & \\
$d_{\min}$ & kódtávolság, két kódszó legkisebb eltérése & \\
$B^T$ & a $B$ paritásrész transzponáltja & \\
$s$ & szindróma & \\
$e$ & a becsült hibavektor: ahol 1-es, ott romlott el a bit & \\
$\hat c$ & a javított (becsült) kódszó & \\
$u_i$ & az üzenet i-edik bitje & \\
$x$ & tetszőleges bit & \\
\end{jeltabla}
\end{jeloles}

\lepes{A generátormátrix}
A $G$ sorai az 1000, 0100, 0010, 0001 üzenetek kódszavai. Például $u=1000$: $p_1=1$, $p_2=1$, $p_3=0$.

\begin{center}\adjustbox{max width=\linewidth}{$\displaystyle 1000\to1000\,110\quad0100\to0100\,101\quad0010\to0010\,011\quad0001\to0001\,111$}\end{center}
\jegyzet{behelyettesítés}{a paritásképletekből}

\begin{center}\adjustbox{max width=\linewidth}{$\displaystyle G=[I\,|\,B],\qquad B=\begin{bmatrix}1&1&0\\1&0&1\\0&1&1\\1&1&1\end{bmatrix}$}\end{center}
\jegyzet{definíció}{}

\lepes{a) A kódszavak}
Minden kódszó a $G$ azon sorainak XOR-ja, ahol az üzenetben 1-es áll.

\begin{center}\adjustbox{max width=\linewidth}{$\displaystyle 0000000,\ 0001111,\ 0010011,\ 0011100,\ 0100101,\ 0101010,\ 0110110,\ 0111001$}\end{center}
\jegyzet{számérték}{}

\begin{center}\adjustbox{max width=\linewidth}{$\displaystyle 1000110,\ 1001001,\ 1010101,\ 1011010,\ 1100011,\ 1101100,\ 1110000,\ 1111111$}\end{center}
\jegyzet{számérték}{a legkisebb nem nulla súly 3, $d_{\min}=3$}

\lepes{b) Paritásellenőrző mátrix}
$H^T$ felül $B$, alul a $3\times3$-as egységmátrix. A szindróma annyi bites, ahány paritásbit van.

\begin{center}\adjustbox{max width=\linewidth}{$\displaystyle H^T=\begin{bmatrix}1&1&0\\1&0&1\\0&1&1\\1&1&1\\1&0&0\\0&1&0\\0&0&1\end{bmatrix}$}\end{center}
\jegyzet{definíció}{3 bites szindróma}

\begin{center}\adjustbox{max width=\linewidth}{$\displaystyle H=[B^T\,|\,I]=\begin{bmatrix}1&1&0&1&1&0&0\\1&0&1&1&0&1&0\\0&1&1&1&0&0&1\end{bmatrix}$}\end{center}
\jegyzet{definíció}{$B$ transzponáltja mellé a $3\times3$-as egységmátrix}

\lepes{c) Dekódolás}
${\cc{v}}={\cc{0110000}}$ a 2. és a 3. helyen 1-es, ezért $H^T$ 2. és 3. sorát adjuk össze.

\begin{center}\adjustbox{max width=\linewidth}{$\displaystyle s=101\oplus011=110$}\end{center}
\jegyzet{behelyettesítés}{}

\begin{center}\adjustbox{max width=\linewidth}{$\displaystyle 110=H^T\ \text{1. sora}\ \Rightarrow\ e=1000000$}\end{center}
\jegyzet{ellenőrzés}{egy hibát feltételezve az 1. bit romlott el}

\begin{center}\adjustbox{max width=\linewidth}{$\displaystyle \hat c={\cc{0110000}}\oplus1000000=1110000\ \Rightarrow\ \hat u=1110$}\end{center}
\jegyzet{számérték}{}

\begin{center}\adjustbox{max width=\linewidth}{$\displaystyle d_{\min}=3$}\end{center}
\jegyzet{ellenőrzés}{biztos csak akkor, ha legfeljebb 1 hiba történt}

\lepes{d) A második kód ($C_2$) generátormátrixa}
Kifejtjük a negyedik paritásbitet: minden $u_i$ annyiszor szerepel, ahány képletben benne van.

\begin{center}\adjustbox{max width=\linewidth}{$\displaystyle p_4=p_1\oplus p_2\oplus p_3=u_1\oplus u_1\oplus u_2\oplus u_2\oplus u_3\oplus u_3\oplus u_4\oplus u_4\oplus u_4=u_4$}\end{center}
\jegyzet{átalakítás}{$x\oplus x=0$}

\begin{center}\adjustbox{max width=\linewidth}{$\displaystyle G_2=\begin{bmatrix}1&0&0&0&1&1&0&0\\0&1&0&0&1&0&1&0\\0&0&1&0&0&1&1&0\\0&0&0&1&1&1&1&1\end{bmatrix}$}\end{center}
\jegyzet{definíció}{az 1. sor súlya 3, ezért $d_{\min}$ marad 3}
\begin{eredmeny}

\begin{center}\adjustbox{max width=\linewidth}{$\displaystyle \text{b) 3 bites szindróma}\qquad\text{c) }\hat u=1110,\ \text{nem biztos (csak egy hibánál)}\qquad\text{d) }p_4=u_4,\ d_{\min}(C_2)=3$}\end{center}
\end{eredmeny}
\begin{tipp}[Ellenőrzés]
\begin{itemize}
\item $C_1$ a (7,4)-es Hamming-kód: $H^T$ hét sora a hét különböző nem nulla 3 bites szó, ezért minden egyes hiba más szindrómát ad.
\end{itemize}
\end{tipp}


% ------------------------------------------------------------------
\section{\texorpdfstring{6.2. feladat: A, B, C forrás Huffman-kódja}{6.2. feladat: A, B, C forrás Huffman-kódja}}\label{o6_02}
\setcounter{lepes}{0}
\begin{feladat}{6.2. feladat}
20 000 szimbólum/s, valószínűségek 50\%, 25\%, 25\%. Huffman-kód, várható bitsebesség, és mennyire közelíti a veszteségmentes kódolás alsó határát?

{\footnotesize\sffamily Forrás: 5-ösöm 18. oldal $\cdot$ 5o\_18\_1. Címkék: vizsga}
\end{feladat}
\begin{keresunk}
\begin{itemize}
\item a Huffman-kód, $\bar L$ (átlagos kódszóhossz), a bitsebesség, és $\bar L$ összevetése $H$-val (entrópia)
\end{itemize}
\end{keresunk}
\begin{jeloles}
\begin{jeltabla}
$\ca{p_A}$ & az A szimbólum valószínűsége & $0{,}5$\\
$\cb{p_B}$ & a B szimbólum valószínűsége & $0{,}25$\\
$\cc{p_C}$ & a C szimbólum valószínűsége & $0{,}25$\\
$\cd{r}$ & szimbólumsebesség (szimbólum/s) & $20$\,s\\
$\bar L$ & átlagos kódszóhossz & \emph{keresett}\\
$A$ & az A szimbólum & \\
$B$ & a B szimbólum & \\
$C$ & a C szimbólum & \\
\end{jeltabla}
\end{jeloles}

\lepes{Huffman-kód}
A két legkisebbet (B és C, 0,25 és 0,25) összevonjuk: 0,5. Marad A (0,5) és BC (0,5). Az ágakra 0-t és 1-et írunk.

\begin{center}\adjustbox{max width=\linewidth}{$\displaystyle A\to0\qquad B\to10\qquad C\to11$}\end{center}
\jegyzet{definíció}{}

\lepes{Átlagos hossz és bitsebesség}


\begin{center}\adjustbox{max width=\linewidth}{$\displaystyle \bar L={\ca{0{,}5}}\cdot1+0{,}25\cdot2+0{,}25\cdot2=1{,}5\ \text{bit/szimbólum}$}\end{center}
\jegyzet{definíció}{}

\begin{center}\adjustbox{max width=\linewidth}{$\displaystyle 1{,}5\cdot{\cd{20}}\,000=30\,000\ \text{bit/s}$}\end{center}
\jegyzet{számérték}{}

\lepes{Az alsó határ}
A veszteségmentes kódolás alsó határa az entrópia.

\begin{center}\adjustbox{max width=\linewidth}{$\displaystyle H={\ca{0{,}5}}\log_22+2\cdot0{,}25\log_24={\ca{0{,}5}}+1=1{,}5\ \text{bit}$}\end{center}
\jegyzet{definíció}{}

\begin{center}\adjustbox{max width=\linewidth}{$\displaystyle \bar L=H$}\end{center}
\jegyzet{ellenőrzés}{pontosan eléri, mert minden valószínűség kettő hatványa}
\begin{eredmeny}

\begin{center}\adjustbox{max width=\linewidth}{$\displaystyle A\to0,\ B\to10,\ C\to11\qquad30\,000\ \text{bit/s}\qquad\bar L=H=1{,}5\ \text{bit: eléri az alsó határt}$}\end{center}
\end{eredmeny}


% ------------------------------------------------------------------
\section{\texorpdfstring{6.3. feladat: Kétbites üzenetek: G, H alakban, szindróma}{6.3. feladat: Kétbites üzenetek: G, H alakban, szindróma}}\label{o6_03}
\setcounter{lepes}{0}
\begin{feladat}{6.3. feladat}
00 $\to$ 00000, 01 $\to$ 01110, 10 $\to$ 10101, 11 $\to$ 11011. a) Szisztematikus G, legvalószínűbb kódszó $v=01010$-ra. b) H és a szindróma.

{\footnotesize\sffamily Forrás: 5-ösöm 19. oldal $\cdot$ 5o\_19\_1. Címkék: vizsga}
\end{feladat}
\begin{keresunk}
\begin{itemize}
\item a) $G$ (generátormátrix) és $\hat c$ (a javított (becsült) kódszó)
\item b) $H$ (paritásellenőrző mátrix) és $s$ (szindróma)
\end{itemize}
\end{keresunk}
\begin{jeloles}
\begin{jeltabla}
$\ca{k}$ & az üzenet hossza bitben & $2$\\
$\cb{n}$ & a kódszó hossza bitben & $5$\\
$\cc{v}$ & a vett szó & $01010$\\
$G$ & generátormátrix & \emph{keresett}\\
$\hat c$ & a javított (becsült) kódszó & \emph{keresett}\\
$s$ & szindróma & \emph{keresett}\\
$I$ & egységmátrix & \\
$B$ & a generátormátrix paritásrésze & \\
$d$ & két szó Hamming-távolsága: hány helyen különböznek & \\
$\hat u$ & a becsült üzenet & \\
$B^T$ & a $B$ paritásrész transzponáltja & \\
$H^T$ & a paritásellenőrző mátrix transzponáltja & \\
\end{jeltabla}
\end{jeloles}

\lepes{a) Generátormátrix}
A $G$ sorai a 10 és a 01 üzenet kódszavai, így az eleje egységmátrix.

\begin{center}\adjustbox{max width=\linewidth}{$\displaystyle G=\begin{bmatrix}1&0&1&0&1\\0&1&1&1&0\end{bmatrix}=[I\,|\,B],\quad B=\begin{bmatrix}1&0&1\\1&1&0\end{bmatrix}$}\end{center}
\jegyzet{definíció}{}

\lepes{a) A legközelebbi kódszó}
Megszámoljuk, hány bitben tér el ${\cc{v}}$ az egyes kódszavaktól.

\begin{center}\adjustbox{max width=\linewidth}{$\displaystyle d({\cc{v}},00000)=2\quad d({\cc{v}},01110)=1\quad d({\cc{v}},10101)=5\quad d({\cc{v}},11011)=2$}\end{center}
\jegyzet{számérték}{}

\begin{center}\adjustbox{max width=\linewidth}{$\displaystyle \hat c=01110\ \Rightarrow\ \hat u=01$}\end{center}
\jegyzet{számérték}{egyetlen bitben tér el}

\lepes{b) H és a szindróma}
$H=[B^T\,|\,I]$. A szindróma $H^T$ azon sorainak XOR-ja, ahol ${\cc{v}}$-ben 1-es áll (2. és 4. hely).

\begin{center}\adjustbox{max width=\linewidth}{$\displaystyle H=\begin{bmatrix}1&1&1&0&0\\0&1&0&1&0\\1&0&0&0&1\end{bmatrix}$}\end{center}
\jegyzet{definíció}{}

\begin{center}\adjustbox{max width=\linewidth}{$\displaystyle s=110\oplus010=100$}\end{center}
\jegyzet{behelyettesítés}{$H^T$ 2. sora 110, 4. sora 010}

\begin{center}\adjustbox{max width=\linewidth}{$\displaystyle 100=H^T\ \text{3. sora}\ \Rightarrow\ \hat c={\cc{01010}}\oplus00100=01110$}\end{center}
\jegyzet{ellenőrzés}{egyezik az a) résszel}
\begin{eredmeny}

\begin{center}\adjustbox{max width=\linewidth}{$\displaystyle \text{a) }\hat c=01110\ (\hat u=01)\qquad\text{b) }s=100,\ \text{a 3. bit hibás}$}\end{center}
\end{eredmeny}


% ------------------------------------------------------------------
\section{\texorpdfstring{6.4. feladat: $G=[10101;\,01011]$}{6.4. feladat: G=[10101;01011]}}\label{o6_04}
\setcounter{lepes}{0}
\begin{feladat}{6.4. feladat}
$G=\begin{bmatrix}1&0&1&0&1\\0&1&0&1&1\end{bmatrix}$. a) Kódszavak, legvalószínűbb üzenet $v=11101$-re. b) H és a szindróma. c) Biztosan jelezhető és javítható hibák?

{\footnotesize\sffamily Forrás: 5-ösöm 19. oldal $\cdot$ 5o\_19\_2. Címkék: vizsga}
\end{feladat}
\begin{keresunk}
\begin{itemize}
\item a) a kódszavak és $\hat u$ (a becsült üzenet)
\item b) $H$ (paritásellenőrző mátrix) és $s$ (szindróma)
\item c) jelezhető és javítható hibák száma
\end{itemize}
\end{keresunk}
\begin{jeloles}
\begin{jeltabla}
$G$ & generátormátrix & \\
$B$ & a generátormátrix paritásrésze & \\
$\ca{v}$ & a vett szó & $11101$\\
$\hat u$ & a becsült üzenet & \emph{keresett}\\
$s$ & szindróma & \emph{keresett}\\
$d$ & két szó Hamming-távolsága: hány helyen különböznek & \\
$\hat c$ & a javított (becsült) kódszó & \\
$B^T$ & a $B$ paritásrész transzponáltja & \\
$H^T$ & a paritásellenőrző mátrix transzponáltja & \\
$d_{\min}$ & kódtávolság, két kódszó legkisebb eltérése & \\
\end{jeltabla}
\end{jeloles}

\lepes{a) Kódszavak}
A négy üzenet: $00\to0$, $10\to$ 1. sor, $01\to$ 2. sor, $11\to$ a két sor XOR-ja.

\begin{center}\adjustbox{max width=\linewidth}{$\displaystyle 00000,\quad10101,\quad01011,\quad10101\oplus01011=11110$}\end{center}
\jegyzet{számérték}{}

\lepes{a) A legközelebbi kódszó}


\begin{center}\adjustbox{max width=\linewidth}{$\displaystyle d({\ca{v}},00000)=4\quad d({\ca{v}},10101)=1\quad d({\ca{v}},01011)=3\quad d({\ca{v}},11110)=2$}\end{center}
\jegyzet{számérték}{}

\begin{center}\adjustbox{max width=\linewidth}{$\displaystyle \hat c=10101\ \Rightarrow\ \hat u=10$}\end{center}
\jegyzet{számérték}{}

\lepes{b) H és a szindróma}
$H=[B^T\,|\,I]$, $H^T$ sorai: 101, 011, 100, 010, 001. A ${\ca{v}}$ egyesei az 1., 2., 3. és 5. helyen vannak.

\begin{center}\adjustbox{max width=\linewidth}{$\displaystyle H=\begin{bmatrix}1&0&1&0&0\\0&1&0&1&0\\1&1&0&0&1\end{bmatrix}$}\end{center}
\jegyzet{definíció}{}

\begin{center}\adjustbox{max width=\linewidth}{$\displaystyle s=101\oplus011\oplus100\oplus001=011$}\end{center}
\jegyzet{behelyettesítés}{}

\begin{center}\adjustbox{max width=\linewidth}{$\displaystyle 011=H^T\ \text{2. sora}\ \Rightarrow\ \hat c={\ca{11101}}\oplus01000=10101$}\end{center}
\jegyzet{ellenőrzés}{egyezik az a) résszel}

\lepes{c) Hibajelzés és -javítás}
A nem nulla kódszavak súlya 3, 3, 4, ezért $d_{\min}=3$.

\begin{center}\adjustbox{max width=\linewidth}{$\displaystyle d_{\min}-1=2\ \text{jelezhető}\qquad\left\lfloor\tfrac{d_{\min}-1}{2}\right\rfloor=1\ \text{javítható}$}\end{center}
\jegyzet{tétel}{}
\begin{eredmeny}

\begin{center}\adjustbox{max width=\linewidth}{$\displaystyle \text{a) }\hat u=10\qquad\text{b) }s=011\qquad\text{c) 2 hibát jelez, 1-et javít biztosan}$}\end{center}
\end{eredmeny}


% ------------------------------------------------------------------
\section{\texorpdfstring{6.5. feladat: Szemcsék a tartályban: a legrövidebb kód}{6.5. feladat: Szemcsék a tartályban: a legrövidebb kód}}\label{o6_05}
\setcounter{lepes}{0}
\begin{feladat}{6.5. feladat}
Hártyával két félre osztott tartályban kék és piros szemcsék mozognak véletlenszerűen. Óránként azt kell rögzíteni, melyik félben milyen színű szemcsék vannak (a számuk nem lényeges). a) Alsó korlát a legrövidebb átlagos kódszóhosszra? b) Szerkessze meg a kódot!

{\footnotesize\sffamily Forrás: 5-ösöm 20. oldal $\cdot$ 5o\_20\_1. Címkék: vizsga, javítás, nehezebb}
\end{feladat}
\begin{keresunk}
\begin{itemize}
\item a) $H$: az alsó korlát
\item b) egy kód és $\bar L$ (átlagos kódszóhossz)
\end{itemize}
\end{keresunk}
\begin{jeloles}
\begin{jeltabla}
$\bar L$ & átlagos kódszóhossz & \emph{keresett}\\
\end{jeltabla}
\par\smallskip\textbf{Adott még:}
\begin{itemize}
\item két szín, két fél, a szemcsék száma nem számít
\item a valószínűségek nincsenek megadva: egyenlőnek vesszük őket
\end{itemize}
\end{jeloles}
\begin{elmelet}[Miért így? Egyszerűen]
\textbf{Miért 9 állapot és nem 16?}\par Mindkét színből van szemcse, így egy szín a bal, a jobb vagy mindkét félben lehet: színenként 3 lehetőség, összesen $3\cdot3=9$. A „sehol sincs” eset kizárt, ezért nem $2^4=16$.
\end{elmelet}

\lepes{a) Alsó korlát}
Kilenc egyenlő valószínűségű állapot entrópiája:

\begin{center}\adjustbox{max width=\linewidth}{$\displaystyle H=\log_29=3{,}17\ \text{bit}$}\end{center}
\jegyzet{definíció}{}

\lepes{b) Huffman-kód}
3 biten $2^3=8$ szó fér el, nekünk 9 kell. Ezért a 8 háromjegyű szó közül egyet kettébontunk két négyjegyűre: 7 szó 3 bites, 2 szó 4 bites. (Egyenlő valószínűségnél a Huffman-eljárás is ezt adja. Kraft: $7/8+2/16=1$.)

\begin{center}\adjustbox{max width=\linewidth}{$\displaystyle \bar L=\frac{7\cdot3+2\cdot4}{9}=\frac{29}{9}=3{,}22\ \text{bit}$}\end{center}
\jegyzet{definíció}{}

\begin{center}\adjustbox{max width=\linewidth}{$\displaystyle 3{,}17\le3{,}22<4$}\end{center}
\jegyzet{ellenőrzés}{több óra adatait együtt kódolva (blokk-kódolással) 3,17-hez tetszőlegesen közel juthatunk}
\begin{tipp}[A gyűjtemény megoldása]
A gyűjtemény szerint az alsó korlát 4 bit, és a nyers 4 bites állapotleírást adja kódnak.
\end{tipp}
\begin{figyelem}[Javítás]
9 lehetséges állapot van, nem 16, ezért 4 bit nem az alsó korlát. Ha a valószínűségek nem egyenlők \texttt{[ELLENŐRIZENDŐ]}, az entrópia kisebb: sok szemcsénél a „mindkét szín mindkét félben” állapot szinte biztos.
\end{figyelem}
\begin{eredmeny}

\begin{center}\adjustbox{max width=\linewidth}{$\displaystyle \text{a) }\bar L\ge\log_29=3{,}17\ \text{bit}\qquad\text{b) Huffman: 7 db 3 bites és 2 db 4 bites szó},\ \bar L=3{,}22\ \text{bit}$}\end{center}
\end{eredmeny}


% ------------------------------------------------------------------
\section{\texorpdfstring{6.6. feladat: Kraft: 2,2,3,3,3,4,4,4,4 és 2,2,3,3,3,5,5,5,5}{6.6. feladat: Kraft: 2,2,3,3,3,4,4,4,4 és 2,2,3,3,3,5,5,5,5}}\label{o6_06}
\setcounter{lepes}{0}
\begin{feladat}{6.6. feladat}
a) Lehet-e egyértelműen megfejthető kódot készíteni a 2,2,3,3,3,4,4,4,4 hosszakkal? b) Prefixmentes kód a 2,2,3,3,3,5,5,5,5 hosszakkal.

{\footnotesize\sffamily Forrás: 5-ösöm 20. oldal $\cdot$ 5o\_20\_2. Címkék: vizsga}
\end{feladat}
\begin{keresunk}
\begin{itemize}
\item a) létezik-e egyértelműen megfejthető kód
\item b) egy prefixmentes kód
\end{itemize}
\end{keresunk}
\begin{jeloles}
\begin{jeltabla}
$l_i$ & az i-edik kódszó hossza & \\
\end{jeltabla}
\par\smallskip\textbf{Adott még:}
\begin{itemize}
\item $l_i$: a kódszavak hosszai
\end{itemize}
\end{jeloles}

\lepes{a) Kraft-összeg}
A McMillan-tétel szerint egyértelműen megfejthető kód csak akkor létezik, ha $\sum2^{-l_i}\le1$.

\begin{center}\adjustbox{max width=\linewidth}{$\displaystyle 2\cdot\tfrac14+3\cdot\tfrac18+4\cdot\tfrac1{16}=0{,}5+0{,}375+0{,}25=1{,}125>1$}\end{center}
\jegyzet{behelyettesítés}{nem létezik}

\lepes{b) Kraft-összeg}


\begin{center}\adjustbox{max width=\linewidth}{$\displaystyle 2\cdot\tfrac14+3\cdot\tfrac18+4\cdot\tfrac1{32}=0{,}5+0{,}375+0{,}125=1$}\end{center}
\jegyzet{behelyettesítés}{létezik, és a kódfa teljes}

\lepes{b) A kód megszerkesztése}
Kanonikus módszer: a hosszak növekvő sorrendjében mindig az előző kódszó + 1-et vesszük, és ha nő a hossz, jobbról 0-kkal egészítjük ki.

\begin{center}\adjustbox{max width=\linewidth}{$\displaystyle 00,\ 01\ \to\ 100,\ 101,\ 110\ \to\ 11100,\ 11101,\ 11110,\ 11111$}\end{center}
\jegyzet{számérték}{$01+1=10\to100$, $110+1=111\to11100$}
\begin{eredmeny}

\begin{center}\adjustbox{max width=\linewidth}{$\displaystyle \text{a) nem (Kraft-összeg }1{,}125>1)\qquad\text{b) }00,\ 01,\ 100,\ 101,\ 110,\ 11100,\ 11101,\ 11110,\ 11111$}\end{center}
\end{eredmeny}
